\chapter{Einbindung in das Gesamtsystem}
\label{ch:integration}

Modul~4 wurde zunächst in einem eigenen Repository entwickelt und anschließend in das
Repository der Gruppe übernommen \quelle{INTEGRATION.md}. Dort ist es auf zwei Wegen
erreichbar: als Greifvariante im automatischen Ablauf und als eigener Reiter in der
gemeinsamen Bedienoberfläche.

% ==============================================================================
\section{Schnittstelle zum Ablauf}
\label{sec:grasper}

Der Ablauf des Gesamtsystems ist ein Zustandsautomat, in dem das Greifen über eine
Schnittstelle \code{Grasper} angesprochen wird. Die Variante von Modul~4 übernimmt nur den
Zustand PICK; das Ablegen auf der Prüfvorrichtung macht weiterhin Modul~2.

\paragraph{Eigene Python-Umgebung.} LeRobot~0.6.1 setzt Python~3.12 voraus, und seine
Abhängigkeiten kollidieren mit denen der Texterkennung von Modul~3. Die Policy läuft deshalb
als Unterprozess in einer eigenen Umgebung. Die Bedienoberfläche und die übrigen Module
bleiben in ihrer Umgebung. Der Nachteil ist eine Startzeit von einigen Sekunden je Fahrt; der
Vorteil, dass keine der beiden Umgebungen auf Kompromissversionen festgelegt werden muss.

\paragraph{Ablauf eines Greifvorgangs} (\cref{fig:grasper}):
\begin{enumerate}
  \item Modul~1 gibt die Draufsichtkamera frei, damit Modul~4 sie öffnen kann; Modul~2 trennt
    seine Verbindung zur Steuerung. Je Kamera und je Steuerung darf immer nur ein Prozess
    verbunden sein.
  \item Der Unterprozess startet die Policy und fährt den Roboter. Die Lage aus der
    Objekterkennung wird übergeben, aber nur protokolliert; die Policy findet das Bauteil
    selbst in den Bildern.
  \item Der Fortschritt wird aus der Ausgabe des Unterprozesses gelesen. Erreicht die Policy
    den Übergabepunkt nicht oder dauert die Fahrt länger als \SI{300}{\second}, gilt das
    Greifen als fehlgeschlagen, und der Ablauf geht in den Fehlerzustand.
  \item Für das Ablegen verbindet sich Modul~2 wieder, setzt Werkzeug, Override und
    Bewegungsprogramm auf seine eigenen Werte zurück und legt das Bauteil ab.
\end{enumerate}

\begin{figure}[htbp]
  \centering
  \begin{tikzpicture}[font=\small, x=1cm, y=1cm]
    \foreach \x/\n in {0/Ablauf, 3.4/Modul~1, 6.8/{Modul~4\\(Unterprozess)}, 10.2/Modul~2} {
      \node[draw, rounded corners=1.5mm, fill=black!5, text width=24mm, align=center,
            minimum height=9mm, font=\scriptsize] at (\x,0) {\n};
      \draw[black!30, densely dashed] (\x,-0.45) -- (\x,-6.3);
    }
    \def\msg#1#2#3#4{\draw[-{Stealth[length=2mm]}, thick] (#1,#3) -- node[above, font=\scriptsize, fill=white, inner sep=1pt] {#4} (#2,#3);}
    \msg{0}{3.4}{-1.0}{Lage bestimmen, Kamera freigeben}
    \msg{0}{10.2}{-1.7}{Verbindung trennen}
    \msg{0}{6.8}{-2.4}{Policy starten (Lage nur Protokoll)}
    \fill[blue!10] (6.55,-2.6) rectangle (7.05,-4.4);
    \node[font=\scriptsize, anchor=west, align=left] at (7.15,-3.5) {greifen, fahren\\bis PRE\_PLACE};
    \draw[-{Stealth[length=2mm]}, thick, densely dashed] (6.8,-4.6) -- node[above, font=\scriptsize] {Übergabepunkt erreicht} (0,-4.6);
    \msg{0}{10.2}{-5.4}{verbinden, Einstellungen zurücksetzen, ablegen}
  \end{tikzpicture}
  \caption[Ablauf eines Greifvorgangs mit Modul~4]{Ablauf eines Greifvorgangs mit Modul~4
  im Gesamtsystem (vereinfacht).}
  \label{fig:grasper}
\end{figure}

\todo[inline]{Höhe des Übergabepunkts über der Prüfvorrichtung vereinheitlichen: Hindernismodell
nennt \SI{7}{\centi\metre}, Orchestrator etwa \SI{9}{\centi\metre}.}

% ==============================================================================
\section{Bedienoberfläche}
\label{sec:gui}

Die Bedienoberfläche von Modul~4 ist als Reiter in die gemeinsame Oberfläche eingebunden und
lässt sich auch eigenständig starten; es ist dieselbe Oberfläche. Sie gliedert sich in vier
Bereiche:
\begin{description}
  \item[Systemcheck] Prüfung von Umgebung, Grafikkarte, Kameras, Verbindung zur Steuerung,
    Kinematik und Sollwerttakt; die Ergebnisse werden als Bericht gespeichert.
  \item[Aufnahme] Ablauftabelle mit Teachen und Touch-up, Tischpunkte, Kameraeinstellungen,
    Einstellungen für Rauschen, Versatz und Objektlagen, Start der Aufnahme und Abdeckungskarte
    der bisherigen Lagen.
  \item[Training] Export und Training mit Fortschrittsanzeige.
  \item[Betrieb] Auswahl eines Modells, Fahrten mit Live-Ansicht, Probegriff bei bekannter
    Lage.
\end{description}
Dass immer nur eine Seite den Roboter und die Draufsichtkamera benutzt, regelt die gemeinsame
Oberfläche. Die Bedienung im Einzelnen beschreibt die Bedienanleitung (Dokument~C).
